\documentclass[11pt]{article}
\usepackage[margin=1in]{geometry}
\usepackage[colorlinks=true,allcolors=blue]{hyperref}
\usepackage{xcolor}

\newenvironment{refquote}%
  {\begin{quote}\itshape\color{black!70}}%
  {\end{quote}}

\setlength{\parskip}{0.5\baselineskip}
\setlength{\parindent}{0pt}

\begin{document}

\begin{center}
{\large\bfseries Response to the Referee}\\[2pt]
Manuscript LR20628MR --- ``Yukawa screening derivation of the bond-valence rule''
\end{center}

We thank the referee for a careful, constructive report.  The
perspective of a condensed-matter reader from outside the bond-valence
community proved especially valuable, and we agree that it matches the
readership of Physical Review Materials.

The referee's paired observations---that the presentation was abrupt,
and that the Letter was short enough to absorb more explanation---%
prompted a more extensive expansion than the specific items below.  We
have enlarged the presentation throughout the manuscript.  The
responses list the specific changes in detail, but additional context
has been added at many points beyond those listed, including a
derivation of the Yukawa flux factor at first use, status statements
distinguishing postulate from approximation from exact identity at
each stage, and mechanistic explanations accompanying the validation
figures.  

The expansion also produced two new results.  First, the
coordination-resolved line family is now closed by explicit intercept
and intersection relations, which give the characteristic softness
$B^*$ a geometric meaning as the logarithmic dilation rate of the
coordination shell with coordination number.  Second, a second-root
analysis of the screening inversion yields a contact-distance law
against Shannon crystal radii, $|\lambda^*_-|=r_{\mathrm{cr}}+c(z)$
with $c(z)\approx 1.05+0.09\,z$~\AA, anchoring the fitted screening
length in a century of independent crystallographic systematics and
identifying a valence-dependent effective oxygen size.  Both results
are validated in the revised Supplemental Material.

\section*{The referee's report, with responses}

\textbf{Block 1 --- introductory remarks.}

\begin{refquote}
``This manuscript presents an interesting and potentially important
attempt to place the bond-valence model on a more fundamental physical
footing by deriving its exponential form as the leading-order limit of
Yukawa-screened electrostatics.  The idea is original, and if correct
it would provide a welcome physical interpretation of empirical
bond-valence parameters that have been widely used for decades.  The
manuscript is also supported by a substantial amount of numerical
analysis and by extensive supplementary material.

I should state, however, that I am not a specialist of the bond-valence
model.  My comments therefore reflect the perspective of a reader
familiar with condensed-matter physics and electronic structure, but
not with the technical details of bond-valence parameterization.  I
think this perspective may actually be relevant for the readership of
Physical Review Materials, since many readers will likely be in a
similar situation.''
\end{refquote}

\textbf{Response.}  We thank the referee for this generous assessment,
and we agree that the stated perspective matches the journal's
readership.

\textbf{Block 2 --- construction of the $B(R_0)$ lines and prior work.}

\begin{refquote}
``My main concern is that the introduction of the theory is rather
abrupt.  Several key objects are introduced without sufficient
explanation for readers outside the bond-valance community.  In
particular, the coordination-resolved $B(R_0)$ lines, which play a
central role throughout the manuscript and in the Supplementary
Information, appear almost without motivation.  From the text, it is
not immediately clear how these lines are constructed numerically.
Are they obtained by fixing $R_0$ and re-optimizing $B$, or by some
other procedure?  Since the subsequent theoretical analysis is built
upon these lines, I think the fitting procedure should be described
explicitly in the manuscript.  In addition, I think the Introduction
would benefit from a broader discussion of previous work aimed at
providing a physical interpretation of the bond-valence model and
related electrostatic approaches, so that the reader can more clearly
appreciate what is genuinely new in the present contribution.''
\end{refquote}

\textbf{Response.}  This block raises three connected requests, and each
are addressed below.

\emph{Construction of the lines.}  The coordination-resolved
$B$--$R_0$ lines were first introduced in the American Mineralogist
paper of Li \emph{et al.}\ (Ref.~[19] of the manuscript), where the
underlying bond-valence solver and the per-structure fits are
developed in detail, and the original manuscript cited that work
intending it to supply this context.  We agree that the construction
was not fully explained and must be stated in the manuscript itself.  
The fitting procedure is now described explicitly in the manuscript, 
with the attribution to Li \emph{et al.}\ made at the point 
of use, and the referee's question is answered directly.  
Both parameters are fitted jointly by linear least squares for each structure, not by
re-optimizing $B$ at fixed $R_0$ (Empirical validation).  Each
DFT-relaxed structure thus contributes one $(R_0,B)$ pair, and the
coordination-resolved line of a species at coordination $n$ is an
outlier-robust regression of $B$ on $R_0$ across those per-structure
pairs (RANSAC consensus refit by least squares on its inliers, at
least five structures per cell), with a fitted slope accepted only
with at least four inlier structures, slope standard error at most 5,
and $R^2\ge 0.3$.  The slopes plotted in Fig.~1 are these regression
slopes, and a pointer to this construction now appears at their first
use in the manuscript (Parameter-free slope).  A new Supplemental
Material paragraph (``Line-fitting procedure'') documents the
complete pipeline, including the coordination binning, the RANSAC
residual threshold, the bootstrap error estimation, and the
structure-count weights.

\emph{Abruptness of the theory introduction.}  The Yukawa flux factor
$e^{-R/\lambda}(1+R/\lambda)$ is now derived by applying Gauss's law
to the screened potential at its first use, with the screening
citations moved there, and the flux-ratio condition is motivated
before it is imposed (Screened-flux derivation).

\emph{Prior work.}  The manuscript's introduction has been rewritten.
It now states concretely what Preiser \emph{et al.}\ established with
point-charge calculations, characterizes Brown's development of the
flux picture into a capacitor-based framework, and distinguishes the
Adams line of work as relating $B$ to bond stiffness rather than to
screening in the surrounding medium.  It also adds the context that
makes the novelty clear.  The near-degeneracy of jointly fitted
$(R_0,B)$ is well documented and has commonly been suppressed by
fixing $B\simeq0.37$~\AA, whereas the present work treats that
correlation itself as physical content and predicts its slope with no
fitted constants.  The opening paragraph now also notes the model's
growing use as a structural descriptor in machine-learning studies,
with three added references.  Referencing of previous work has been
improved throughout, including the explicit attribution of the
$B$--$R_0$ line construction to Li \emph{et al.}\ (Ref.~[19]) at its
point of use.

\textbf{Block 3 --- exact reformulation vs.\ approximation vs.\
interpretation.}

\begin{refquote}
``More generally, I found that several conceptual steps in the
derivation would benefit from additional discussion.  For example, the
partition-function formulation and the definition of the
first-principles screening centroid are mathematically interesting,
but their physical motivation is introduced rather quickly.  Readers
unfamiliar with bond-valence theory may struggle to distinguish what
is an exact reformulation, what is an approximation, and what is a
physical interpretation.''
\end{refquote}

\textbf{Response.}  The manuscript now makes this bookkeeping explicit
at every stage.  Its introduction closes with a roadmap stating that
the screened flux-ratio partition is a physical postulate, that the
isotropic Yukawa weight and its linearization are the two controlled
approximations, and that the partition-thermodynamic relations are
exact identities of the bond-valence equation.  The partition equation
is followed by an exactness statement locating the first mathematical
approximation in the linearization.  The manuscript's Partition
thermodynamics section now opens by motivating the ensemble reading
(it supplies the reference length left open in the derivation and the
shell radius that the first-principles comparison requires), declares
its identities exact with the thermodynamic language as
interpretation, and notes where the screened-flux and bond-valence
forms coincide only to first order.  The screening centroid
$r_{\mathrm{eff}}$ is now introduced in the manuscript
(First-principles confirmation) by its physical role---if bond valence
is captured screened flux, the shell radius resolved by the partition
should track where screening charge accumulates along each bond---%
before its operational definition.

\textbf{Block 4 --- spread near $z=3$ in Fig.~1 and slope
uncertainties.}

\begin{refquote}
``I am also only partially convinced by the empirical validation
presented in Fig.~1.  The overall agreement with the predicted slope
is encouraging, but there remains a substantial spread for some
oxidation states, particularly around $z=3$.  The reported weighted
$R^2$ values are excellent, yet the visual impression is considerably
less compelling.  It would therefore be useful to discuss the origin
of this scatter more explicitly, and, if possible, to indicate the
uncertainty associated with the fitted slopes.''
\end{refquote}

\textbf{Response.}  Both requests are implemented, in the manuscript
and in the Supplemental Material, and Fig.~1 itself has been
redesigned.  The manuscript's Fig.~1 is now a box-and-whisker summary
of the per-species slopes at each formal charge, so the spread the
referee noted is displayed rather than implied, and the full
per-species scatter has moved to the Supplemental Material
(``Per-species slopes at fixed coordination'') with species-resolved
error bars.  A new manuscript section, ``Origin of the near-pole
spread,'' analyzes the scatter.  The pole form itself sets the
pattern.  A species-level perturbation $u$ of the inverse slope
shifts the fitted slope by $-\beta^2 u$, so fixed chemical scatter is
amplified quadratically toward the pole, and across the charge
classes of Fig.~1 the per-class spread indeed grows as
$|\beta|^{1.8}$.  The $z=3$ class at $n=4$ is the clearest case.  The
same cations at $n=6$ match the prediction to a median
$|\Delta\beta|=0.03$, so chemical diversity alone is not the origin.
The largest deviations belong to species for which fourfold oxygen
coordination is rare (chiefly the trivalent lanthanides and Y, median
17 structures per fit), whereas well-populated trivalent species lie
on the curve.  Slope uncertainties are now reported in the manuscript
(Parameter-free slope) and documented in the Supplemental Material
(``Slope standard errors'').  Each fitted slope carries a bootstrap
standard error, with median relative standard error $0.2\%$ across
the 150 slopes and 90th percentile $4\%$, so the visible scatter
reflects real differences among species rather than fit uncertainty.

\textbf{Block 5 --- near-coincidence of the coordination-dependent
lines.}

\begin{refquote}
``Conversely, while examining the Supplementary Information, I was
somewhat surprised by a different feature: for many species the
coordination-dependent $B(R_0)$ lines are almost coincident, making
the determination of their intersection appear rather delicate.  This
may simply reflect the logarithmic dependence of the predicted slopes
on $z/n$, so that the expected differences between coordination
numbers are intrinsically small, but if so I think it would be
worthwhile to comment on this point explicitly.  As presented, the
figures leave the impression that the characteristic intersection is
in some cases only weakly constrained.''
\end{refquote}

\textbf{Response.}  The referee's surmise is exactly right, and we now
comment on it explicitly in the manuscript (Characteristic
intersections) and in the Supplemental Material (``Near-coincidence
of same-branch lines'').  Subtracting two predicted slopes gives
$\beta^{(n_2)}-\beta^{(n_1)}=\ln(n_2/n_1)/[\ln(z/n_1)\ln(z/n_2)]$,
which is small whenever $z$ is well separated from both coordination
numbers, so same-branch lines are expected to appear nearly
coincident.  The manuscript states this alongside the conditioning
statement.  The intersection remains constrained exactly to the extent
that the data resolve the shell dilation, and the characteristic pair
is never computed from a single pair of lines---the construction pools
all coordination lines of a species with structure-count weights and
reports the residual spread $\sigma_B$ as its conditioning diagnostic.
A new Supplemental Material paragraph quantifies the conditioning.
Across the 85 species with three or more lines the median
$\sigma_B/B^*$ is $5.8\%$ and 62 fall below $10\%$, while the 18
species constrained by only two lines have $\sigma_B=0$ by
construction and are flagged as provisional in the table.  The
referee is also correct that some intersections are only weakly
constrained, and the revised documents say so rather than leaving it
as an impression.

\textbf{Block 6 --- overall assessment and requested emphases.}

\begin{refquote}
``These comments should not be interpreted as a negative assessment of
the work.  On the contrary, I found the central idea original and
interesting, and I appreciated the transparency of the Supplementary
Information, which documents the fitting procedure and underlying
data in considerable detail.  My impression is that the manuscript
would become considerably more accessible---and, consequently, more
convincing to a broad materials-physics audience---if the authors
devoted somewhat more space to explaining the construction of the
empirical $B(R_0)$ lines, the assumptions entering the derivation,
and the interpretation of the validation figures.

Overall, I believe the manuscript has the potential to make a valuable
contribution, but I would encourage the authors to improve the
presentation and clarify the methodological aspects before
publication.''
\end{refquote}

\textbf{Response.}  We thank the referee for this encouragement.  The
three named emphases are addressed under Blocks~2 and~3 (construction
of the lines; the assumptions entering the derivation) and, for the
interpretation of the validation figures, by two additions.  The
manuscript's Discussion now explains the offset of the centroid
comparison from the identity line mechanistically---the Gibbs shell
centroid references the oxygen nucleus while $r_{\mathrm{eff}}$ tracks
the screening charge displaced into the bond---with the directly
measured oxygen-side shift $\Delta_{\mathrm{O,val}}=0.39\pm0.03$~\AA{}
matching the gap of the fitted lines, and the redesigned Fig.~1
displays the spread that the earlier presentation obscured.

\textbf{Block 7 --- Supplemental captions.}

\begin{refquote}
``Note: many figure captions in the Supplementary seem to be
incorrect.''
\end{refquote}

\textbf{Response.}  We thank the referee for catching this.  A
systematic audit of every caption against the rendered figure pages
confirmed the problem.  The atlas pages order panels alphabetically by
element within each block, while the captions had assumed
atomic-number ordering, so all eight $d$-block captions misidentified
their contents.  Every atlas caption now lists its panels explicitly
and correctly, the ordering convention is stated in the Supplemental
Material's atlas introduction, and the remaining caption inaccuracies
(color versus line-style encodings, species orderings) have been
fixed.  Supplemental figures and tables now carry S-prefixed
numbering to prevent ambiguity with main-text figure references.

\textbf{Block 8 --- Letter format.}

\begin{refquote}
``Justification for/against Letter:

I am not sure if the Letter format is the most appropriate for this
paper.  As now, I find it a bit abrupt, but it is also quite short, so
maybe the authors can add enough explanations while keeping it a
Letter.''
\end{refquote}

\textbf{Response.}  We followed this suggestion.  The explanations
described above were added while condensing elsewhere, and the
revised manuscript stands at approximately 4{,}200 word-equivalents,
within the Letter limit.

\section*{Additional changes}

Beyond the items above, the deeper revision the report prompted led to
the following changes:

\begin{itemize}
\item The manuscript's screening-collapse discussion at $z=n$ was
corrected.  The original text overstated the constraint (the
valence-sum rule fixes only the mean bond valence, not each bond),
and the manuscript now distinguishes truly degenerate shells, for
which a single structure leaves $B$ unconstrained, from nearly
degenerate shells, which are satisfied at finite $B$, with corpus
counts for both.  Relatedly, the branch change across the pole is now
correctly attributed to the sign reversal of the slope $\beta$ rather
than to the sign of $B$.
\item The tabulated screening length $\lambda^*$ now covers both
branches (signed values for the two anti-screening species), and the
second root $\lambda^*_-$ is tabulated for all species, supporting the
new contact-distance analysis.
\item A sign census of the per-shell fits, its band-gap resolution
(Supplemental Fig.~S1), a comparison of $B^*$ with published softness
parameterizations, and the slope-error and line-conditioning
statistics were added to the Supplemental Material, and all quoted
statistics are reproduced by scripts in the project repository.
\item Fig.~2 was redrawn for readability (smaller markers, labels
adjacent to their points, Group-1 labels below and Group-2 labels
above the fit line), and terminology and notation were unified (Gibbs
ensemble language, reference length $R'$ distinguished from the
characteristic intercept $R_0^*$).
\end{itemize}

We believe the revised manuscript is substantially more accessible, and
we thank the referee for a report that improved it well beyond the
specific points raised.

\end{document}
